\chapter{Introduction}
\label{ch:intro}

Books are easier to maintain when each chapter lives in its own file.
Use \verb|\include{filename}| in \texttt{main.tex} (without the \texttt{.tex} extension).

\section{Why this layout?}
The \texttt{book} class supports parts, chapters, and front/back matter.
The table of contents and chapter breaks are generated automatically.

\section{Next steps}
For background on multi-file workflows, see \citep{exampleArticle2024}.
\begin{itemize}
  \item Rename chapters and add new \texttt{.tex} files under the project.
  \item Add figures with \verb|\includegraphics| (upload images via \textbf{+ File}).
  \item Cite sources with \verb|\citep{example2024}| once \texttt{refs.bib} is filled in.
\end{itemize}
